\section{Calibrated models: fluid and discrete-event}
\label{sec:models}

\subsection{Construction}
\label{sec:construction}

Two model layers share one calibration. The fluid model evolves a
session reservoir with completion-coupled turn release, per-replica
decode interference, and a contended fleet restore channel; restores
form a third $\max()$ term in the waiting-time expression, achieved
restore traffic advances the HBM write clock, and the tier churns at
the measured ingest rate with hit-refresh \evid{R9}. Arriving (salted)
sessions enqueue a forced full-miss first request as an explicit
queue class, and surge populations are cancelled at surge end,
matching the protocol.

The DES adds the discrete structure the fluid cannot carry:
watermark admission (waiting requests join a replica's batch, full
KV need reserved, while the batch fits under the occupancy
watermark; the prefill engine stays serialized), a two-clock LRU
(the HBM clock advances with prefill, restores, and generation; the
tier clock with new writes only; touch stamps refreshed at decode
end), and a per-replica FCFS restore channel at $\Br$ with
asynchronous onboarding \evid{R9}. A single-clock tier variant was
rejected against measurement: it drives the tier to exhaustion and a
cold collapse that the 300-min hardware run refutes, and it
contradicts the measured tier-ingest rate.

The deployed-router layer replaces the idealized lexicographic
router with the deployed EPP algorithm, every scheduling constant
read from deployed configuration or plugin code, none fitted:
prefix-affinity sticky threshold 0.80; TTFT load gate 18000\,ms with
TTFT estimated as in-flight tokens over the deployed peak prefill
throughput of 28888 tokens/s; token-load scoring; and the 300\,s
in-flight staleness reap, under which deep engine queues undercount
and the gate rarely triggers under backlog \evid{R9}. The model contains
no coordinator term: an earlier coordinator-saturation freeze is
rejected on two independent grounds - the router-side measurements
of Section~\ref{sec:coordinator} contradict its physical interpretation, and it is
unnecessary (no outcome classification changes without it at the
300-min horizon). The idealized router serves as a negative control:
it produces identical collapse rates at capacity-matched points
(3/5 at both ($N{=}8$, 0.020) and ($N{=}16$, 0.040)), so the measured span
dependence is not reproducible from scheduling idealizations plus
per-replica physics.

\subsection{Calibration}
\label{sec:calibration}

Table~\ref{tab:calibration} lists every constant and its provenance; the single fitted
quantity is the effective tier capacity (an assumed size-linear
nominal times the fitted efficiency).

\begin{table*}[t]
\centering
\small
\begin{tabularx}{\textwidth}{@{}p{0.16\textwidth}p{0.22\textwidth}X@{}}
\toprule
Constant & Value & Source \\
\midrule
Prefill rate $p$ & 17k tokens/s/replica & measured saturated uncached plateaus, 131--138k tokens/s fleet \evid{R9} \\
HBM pool $C$ & 6486 blocks $\times$ 256 tokens per TP4 replica & engine configuration; byte size cross-checked against KV/token \evid{R9} \\
Decode interference & $\mathrm{tpot} = 10.8\,\mathrm{ms} + 2.56 \times 10^{-4} (R/N)^2$ & fit over 552 measured windows (Figure~\ref{fig:tpot}); per-replica form selected, fleet-total form falsified on saturated 4-replica windows (RMSE 27.4 vs 18.6\,ms) \evid{R9} \\
Restore rate $\Br$ & 23k tokens/s/replica effective & measured per-replica restore plateau, 21.4--24.0k across four congested runs at both fleet scales \evid{R4}; a congested-state throughput calibration, not a channel ceiling (Section~\ref{sec:falsifiers}) \\
DMA link & 151.5\,GB/s = 1.19M tokens/s/replica & measured from offload counters; $\Br$ runs the link at $\sim$2\% duty, so per-transfer overhead binds \evid{R4} \\
KV per token & 127\,KB (fp8, 62 layers, 8 KV heads, dim 128) & validated against the pool byte size; windowed host-to-device bytes / 127\,KB matches the windowed tier-hit token rate to under 1\% in every tiered run \evid{R9} \\
$\Ereq$ & 118 realized requests/session & 28,366 records / 240 sessions in the 3\,h windows; the corpus mean of 174 inflates demand $\sim$40\% and moves the model boundary off the measured bracket \\
Reuse-gap CDF & corpus mixture, capped at 10.5\,s & the protocol gap cap, an experimental constant \evid{R9} \\
First-request size & $\sim$51.6k tokens & system prompt plus first input; forced full miss per arriving salted session \\
Router constants & 0.80 / 18000\,ms / 28888 tokens/s / 300\,s & read from deployed configuration and plugin sources; none fitted \evid{R9} \\
Tier nominal capacity & 2520 tokens/replica per configured size unit (size 500 = 1.26M) & ASSUMED size-linear mapping anchored at size 500 to the 160\,GB/replica host-RAM headroom at 127\,KB/token; the configured-unit-to-byte semantics are unconfirmed (Section~\ref{sec:limitations}) \\
Tier efficiency & 0.67 (FITTED) & cold-side ordering of the tier-size series: effective(250) = 0.42M must sit at or below the DES cold/congested transition and effective(375) = 0.63M at or above it, the in-model transition spanning (0.42, 0.63)M; the congested-state miss share independently implies an effective window below nominal \evid{R9}. Because the nominal is assumed, the product (0.42/0.63/0.84M at sizes 250/375/500) is in substance a single fitted effective-capacity scale \\
\bottomrule
\end{tabularx}
\caption{Calibration constants and provenance.}
\label{tab:calibration}
\end{table*}

\subsection{Validation}
\label{sec:validation}

\begin{figure*}[t]
\centering
\includegraphics[width=\textwidth]{overlay_b1_dynamics.png}
\caption{Fluid overlay on the four protocol configurations (0.022,
0.017, 0.012 sessions/s untiered at $N{=}8$; 0.011 with tier at $N{=}4$);
surge window shaded; model curves labeled as model output. The
waiting-queue trajectories match in all four configurations,
including the post-cancellation re-growth in the collapsing one.}
\label{fig:fluid}
\end{figure*}

\begin{figure*}[t]
\centering
\includegraphics[width=\textwidth]{des_b1_overlay.png}
\caption{DES validation configurations, 5 seeds each, min-max
shading; deployed-router outcomes: collapse 4/5 at 0.022 (measured
3/3), recovery 5/5 at 0.017 and 0.012, recovery 5/5 for the
4-replica tiered configuration.}
\label{fig:desb1}
\end{figure*}

\begin{figure*}[t]
\centering
\includegraphics[width=\textwidth]{overlay_nseries.png}
\caption{DES fleet-size sweep (12 seeds $\times$ 6 configurations, 300
min, no coordinator term) with measured runs overlaid; the
16-replica collapse curve sits left of the 8-replica curve at
matched normalized rate. Model output; absolute placement carries
the documented reservation bias.}
\label{fig:nseries}
\end{figure*}

\begin{figure}[t]
\centering
\includegraphics[width=\columnwidth]{fig_hysteresis.png}
\caption{Model coexistence band (DES, deployed router, $N{=}8$,
untiered, 3 seeds per point, 300-min horizon; model output).
End-state hit rate versus session rate for the unperturbed protocol
(warm branch) and the standard-surge protocol. Between roughly 0.017
and 0.022 sessions/s the two branches coexist: unperturbed runs stay
warm where perturbed runs end cold. At higher rates the warm branch
itself decays by the delayed mechanism of Section~\ref{sec:horizon}. Placement
carries the documented leftward reservation bias (Section~\ref{sec:biases}).}
\label{fig:hysteresis}
\end{figure}

\begin{figure}[t]
\centering
\includegraphics[width=\columnwidth]{fig_phase_diagram.png}
\caption{Model outcome map over (session rate, effective tier
capacity) under the standard perturbation (DES, deployed router,
$N{=}8$, 3 seeds per cell, 300-min horizon; model output). The untiered
row reproduces the probability band; the smallest tier falls through
to cold at supercritical rates; larger tiers convert collapse into
persistent congestion. The model produces no complete recoveries in
the congested cells - its documented recovery-branch failure
(Section~\ref{sec:failures}) - so the measured complete-recovery outcomes at sizes
375--500 have no model counterpart in this map.}
\label{fig:phasediagram}
\end{figure}

The model reproduces, in order of evidential weight \evid{R9}:

\begin{itemize}
\item The untiered band and outcome classes. Under the deployed-router
  model: collapse at 0.022 in 4/5 seeds (measured 3/3), recovery 5/5
  at 0.017 and 0.012, and the 4-replica tiered recovery 5/5 at its
  180-min horizon. An earlier idealized-router revision additionally
  matched the in-surge peak waiting depths (620--787 modeled against
  675--790 measured on the 8-replica configurations), showed smooth
  20--40\,min hit-rate erosion where the fluid transition is square,
  and produced the post-cancellation warm interlude in every
  collapsing seed (the measured interlude/no-interlude pair sits
  inside the seed spread).
\item The third regime, all four discriminators in 5/5 seeds: restore
  channel pinned near $N \Br$, linear slow queue divergence, high $h$,
  no cold collapse (Figure~\ref{fig:restore}). Quantitative gaps are stated, not
  fitted: model $h$ 0.95 vs measured 0.70--0.85, restore share 1.0 vs
  0.66--0.79, completion rate 1.6 vs 2.5--4.3 requests/s.
\item The size-250 fall-through in direction and persistence boundary.
  This is calibration-constrained rather than independent
  validation - the tier efficiency is fitted to exactly this
  cold-side ordering - so only the trajectory shape and seed
  structure are non-fitted. Under the final calibration the fine
  capacity sweep gives cold 5/6 at effective 0.42M while the
  size-500 configuration runs congested 6/6, never cold, with the
  DES mean tracking the measured decay-through-congestion
  trajectory (Figure~\ref{fig:tiercap}).
\item The span shift, emergent in direction. With no $N$-dependent term,
  DES collapse counts at 300\,min (12 seeds per configuration) are
  2/12, 11/12, 12/12 at ($N{=}8$) 0.017/0.020/0.022 and 5/12, 12/12,
  11/12 at ($N{=}16$) 0.034/0.037/0.040. The shift is visible at the
  0.017-equivalent pair (5/12 vs 2/12), a seed-level contrast that
  does not reach significance at 12 seeds (one-sided exact $p \sim
  0.19$); the upper matched pair is ceiling-saturated and carries no
  shift information. The direction matches the measured sweep. The
  shift arises from the deployed router's affinity-load arbitration
  coupling every replica's post-cancellation drain race; the
  idealized router - equally fleet-spanning and load-balancing -
  produces no shift (Section~\ref{sec:construction}), so a shared span alone does not
  produce it, and which deployed feature drives it is not isolated
  (Section~\ref{sec:limitations}). The model predicts shard routers behave as
  independent 8-replica systems, matching the measured containment
  \evid{R10}.
\end{itemize}

\subsection{Known biases}
\label{sec:biases}

Full KV need is reserved at batch admission (vLLM allocates blocks
progressively during chunked prefill), which overstates running-set
residency during backlog drains. Both DES collapse curves sit left
of the measured ones as a consequence; the model's quantitative
claim is the shift and the outcome-class structure, never absolute
collapse probabilities \evid{R9}. Further documented limitations: decode
interference is frozen at decode start; there is no
preemption/recompute path; collapsed-state $h$ floors at 0 in most
seeds versus the measured 2--12\% residual.

\subsection{Executed falsifiers}
\label{sec:falsifiers}

Predictions were registered before hardware windows and are reported
regardless of outcome.

\begin{itemize}
\item The 4-replica tiered configuration at 300\,min: predicted waiting
  depth 52--160; measured 1.3, stable 3/3 (one run at the full
  horizon). Falsified; the miss is assigned to the recovery branch
  (Section~\ref{sec:failures}) \evid{R9}.
\item The tier 375/650 bracket, registered under the prior tier
  efficiency of 0.5 (which placed 375 inside the predicted cold
  region): 375 completely recovered 2/2, excluding 0.5 and forcing
  the recalibration to 0.67; the 650 configuration was not
  executable (boot failure, Section~\ref{sec:sizeseries}) \evid{R11}\evid{R9}.
\item The 16-replica edge points: the sweep's left shift is confirmed in
  direction by the measured counts (Section~\ref{sec:leftshift}; the 0.034 recovery
  is $n=1$). The DES collapses 5/12 at 0.034 where the one measured
  run recovered; an $n=1$ measurement is uninformative about that rate
  (a recovery has probability 7/12 under the model's own counts).
\item The fleet-total interference form: registered and falsified by the
  saturated 4-replica windows (predicted 33\,ms vs measured 81\,ms
  mean time-per-token); the per-replica form replaced it
  (Section~\ref{sec:perreplica}) \evid{R6}\evid{R9}.
\item The $\Br$ ceiling interpretation: encoded in the serial FCFS
  channel; falsified by the two complete-recovery runs
  (Section~\ref{sec:pinnedband}). The 23k plateau is a congested-state outcome. The
  serial channel stands as a documented deficiency; its
  concurrent-restore replacement fails the recovery-branch screen
  and is not adopted (Section~\ref{sec:failures}).
\end{itemize}

\subsection{Documented failures and the recovery branch}
\label{sec:failures}

The recovery branch is the model's primary open failure: zero
complete recoveries in 30 sweep runs at ($N{=}8$, 0.022) at any tier
capacity, against 4 measured in 8 runs at sizes 375--500; the missed
complete recovery at the 0.020 boundary; and the falsified 4-replica
prediction above \evid{R9}. The failure is one-sided: the DES reproduces
cold collapse, restore-bound congestion, and the untiered
recoveries, but cannot exit the congested state through drain-back.
The fluid layer carries a second documented failure, branch
selection at cancellation: it reproduces the in-surge restore-bound
state but flushes its restore backlog within $\sim$2\,min of cancellation
and exits to the warm branch at both scales; the DES,
restore-serialized with heterogeneous gaps, diverges - matching
measurement. Both layers sit near the same critical balance, and no
fluid constant flips the branch without violating a measurement, so
the DES is the validated layer for the third regime.

Four candidate mechanisms for the recovery-branch failure were
screened on a shared protocol (seven configurations spanning tier,
untiered, capacity, and span variants; 6 seeds; 300\,min), with the
guard constraint that any candidate preserve the cold/congested
boundary and the untiered band. All four screens are negative; none
is adopted; each is summarized here with its exclusion scope, and
the artifacts carry the full protocols.

\begin{enumerate}
\item Reservation timing (progressive KV allocation, removing the
   full-reservation-at-admission overstatement): negative and
   informative - the screened configuration remains 0/6 recovered,
   and the variant manufactures a non-physical interference-locked
   congested state in an untiered control configuration that the
   engine's step-budget scheduler cannot enter. Its retained value
   is diagnostic: every failing baseline run drains from waiting
   depth 319--745 to a trough of 0--31, then re-enters congestion.
\item Stale-content tier churn (a credit bracket subtracting dead-surge
   write tokens from tier ages - the complete-instantaneous-reclamation
   bound): null; outcome counts match baseline in every
   configuration, and the instrumented run shows live write flux
   churns dead content out of the model's window $\sim$13\,min after
   cancellation with a rescuable live-miss fraction of 0.0
   thereafter. Definitive exclusion within the two-clock
   abstraction.
\item Restore-channel serialization (concurrent onboarding at the
   measured 1.19M tokens/s link with a derived 4.23\,s per-transfer
   overhead): negative; failing configurations still re-enter
   congestion, and the congested pin overshoots the measured
   0.81--1.04 $N \Br$ band (peaks 1.91$\times$).
\item Occupancy-gated HBM eviction (allocation-pressure-only eviction
   replacing the unconditional aging clock, zero new constants,
   pre-quantified as small since the trap itself pins occupancy near
   0.88): negative; control configurations pass, screened
   configurations fail 5/6, and the trap relocates - pending-restore
   entries hold 0.69 of the fleet pool as full-KV reservations
   queued behind the serial channel.
\end{enumerate}

Jointly the screens sharpen the failure to a re-entry trap: failing
model runs drain to near-empty queues, then every completion
re-admits a full-context restore (model restore share 1.0, $\sim$100k
tokens per restore) where the measured congested state runs
$h$ 0.70--0.85 with tier-hit fraction 0.62--0.79. Candidates 3 and 4
bound the trap's carrier - with the channel widened it pins via
elevated restore throughput; with residency corrected it pins via
reservation-holding backlog - so the surviving suspect is the
restored volume per tier hit interacting with full-KV reservation at
admission (partial-context restores, or tier-content staleness
outside the two-clock window model), not channel capacity,
reservation timing, or residency semantics \evid{R9}. We present this as
the diagnosis of an open model failure, not as a result. The
gauge-semantics caveat of Section~\ref{sec:stack} applies to any successor
variant: restore-in-flight entries must not feed the decode
interference term while still holding reservations.
